\section{为什么需要极限}
\label{sec:02-Calculus/02-Limits-and-Continuity/01}

\subsection{测速仪在撒谎吗}

你开车在高速公路上，仪表盘显示 100 km/h。你相信这个数字。

但想一秒钟：``此刻''的速度是什么意思？此刻没有时间流逝，没有距离被走过。速度的公式是``距离除以时间''，而你试图用零长度的距离除以零长度的时间——\(0/0\)，毫无意义。

测速仪不是在撒谎，但它测量的不是``此刻走了多远''。它测量的是：如果你按此刻的状态继续开，一小时会走 100 公里。这是一种\textbf{趋势}判断，不是一次除法。

把这个困境写成数学形式。设你在时刻 \(t\) 的位置为 \(s(t)\)。从 \(t\) 到 \(t+h\) 这个时间窗口内的平均速度：
\[
\frac{s(t+h)-s(t)}{h}.
\]
取具体数字。假设 \(s(t)=4.9t^2\)（自由落体，忽略单位和空气阻力）。在 \(t=1\) 秒处，你想知道瞬时速度。

先算一小段时间内的平均速度：
\[
\begin{aligned}
h=0.1 &: \quad \frac{4.9\times 1.1^2 - 4.9}{0.1} = \frac{5.929-4.9}{0.1} = 10.29,\\[4pt]
h=0.01 &: \quad \frac{4.9\times 1.01^2 - 4.9}{0.01} = \frac{4.99849-4.9}{0.01} = 9.849,\\[4pt]
h=0.001 &: \quad \frac{4.9\times 1.001^2 - 4.9}{0.001} = \frac{4.9098049-4.9}{0.001} = 9.8049.
\end{aligned}
\]

数字在稳定下来——接近 9.8。但你永远不能用 \(h=0\) 直接计算：分母为零，公式炸了。

极限解决的就是这个困境：\textbf{不要求在目标点上有定义，只关心靠近时趋势是否收敛}。

\subsection{代入失败，趋势完全清楚}

上一节我们已经窥见这个思路。再练一个经典例子来建立肌肉记忆：

\[
f(x)=\frac{x^2-1}{x-1}.
\]

\(x=1\) 代入：\(0/0\)，未定义。但 \(x\neq 1\) 时：
\[
\frac{x^2-1}{x-1} = \frac{(x-1)(x+1)}{x-1} = x+1.
\]

所以当 \(x\) 靠近 1 时，\(f(x)\) 靠近 \(1+1=2\)。记为：
\[
\lim_{x\to 1}\frac{x^2-1}{x-1}=2.
\]

这里有一个关键区分：函数值和极限值是两回事。\(f(1)\) 没有定义，极限却是 2。退一步说，即便你强行定义 \(f(1)=100\)，极限仍然是 2——因为极限只看 \(x=1\) \textbf{附近}（但排除 \(x=1\) 自身）的行为。

把函数值和极限分开看待，是微积分能够讨论瞬时变化的前提。瞬时速度、切线斜率、面积密度的精确值——全都依赖这个区分：先在点附近观察趋势，再让观察窗口的宽度趋于零。

\subsection{单侧极限：从左边来和从右边来，未必一致}

有时趋势的方向是故事的核心。从左侧接近 \(a\)（所有小于 \(a\) 的 \(x\) 沿数轴向右逼近）记作 \(x\to a^-\)。从右侧接近（大于 \(a\) 的 \(x\) 向左逼近）记作 \(x\to a^+\)。

符号函数是一个直白的例子：
\[
\operatorname{sgn}(x)=
\begin{cases}
-1, & x<0,\\
 0, & x=0,\\
 1, & x>0.
\end{cases}
\]

从左侧逼近 0：\(x=-0.1, -0.01, -0.001,\ldots\)，函数值始终是 \(-1\)。所以
\[
\lim_{x\to 0^-}\operatorname{sgn}(x) = -1.
\]

从右侧逼近 0：\(x=0.1, 0.01, 0.001,\ldots\)，函数值始终是 1。所以
\[
\lim_{x\to 0^+}\operatorname{sgn}(x) = 1.
\]

两边趋势不同。双侧极限 \(\lim_{x\to 0}\operatorname{sgn}(x)\) 不存在——它不是``存在但算不出来''，而是根本不存在，因为无论你猜哪个数，总有一侧的趋势不指向它。

双侧极限存在当且仅当两个单侧极限都存在且相等：
\[
\lim_{x\to a}f(x) = L
\quad\Longleftrightarrow\quad
\lim_{x\to a^-}f(x) = \lim_{x\to a^+}f(x) = L.
\]

注意 \(f(a)\) 自身的取值与这个判断完全无关。符号函数的 \(f(0)=0\) 没能调和两边的分歧——函数在单点上的定义补救不了左右趋势的割裂。

在定义域的端点处，只有定义域方向允许的单侧极限有意义。\(\sqrt{x}\) 的定义域是 \([0,\infty)\)，\(x=0\) 只有右侧有合法的实数输入，所以讨论 \(x\to 0^-\) 没有意义——那些 \(x\) 不在这道题的论域里。

\subsection{无界增长：当``趋于无穷''是唯一的诚实描述}

看这个：
\[
\lim_{x\to 0}\frac{1}{x^2}.
\]

代入几个数：\(x=0.1\) 得 100，\(x=0.01\) 得 10000，\(x=0.001\) 得 1000000。越小越大，大得没有天花板。你事先随便说一个数——一亿——只要 \(x\) 靠得够近（\(|x|<0.0001\)），函数值就能超过一亿。

记作：
\[
\lim_{x\to 0}\frac{1}{x^2} = +\infty.
\]

\(+\infty\) 不是实数——你不是在说函数``在 0 点取值无穷''。你是在说：\textbf{函数值没有上界，且能持续增长}，对应图像上 \(x=0\) 处有一条竖直渐近线。

但别把符号方向搞混：
\[
\lim_{x\to 0^+}\frac{1}{x} = +\infty,\qquad
\lim_{x\to 0^-}\frac{1}{x} = -\infty.
\]

左右两侧的符号截然相反，所以不能写 \(\lim_{x\to 0}1/x = \infty\) 了事——那会掩盖左侧冲向负无穷的事实。双侧无穷极限要求两侧都趋于同一个带符号的无穷。

\subsection{远处的行为：输入趋于无穷时，输出去哪}

另一种``无穷''问题反过来：不是函数值炸了，而是输入本身越走越远。

\[
\lim_{x\to\infty}\frac{2x+1}{x+3}.
\]

直觉：\(x\) 极大时，分子的 \(2x\) 主导，分母的 \(x\) 主导，比值趋近 \(2/1=2\)。

用数字验证：\(x=100\) 得 201/103 \(\approx\) 1.951，\(x=10000\) 得 20001/10003 \(\approx\) 1.9997。确实在靠近 2。

记为：
\[
\lim_{x\to\infty}\frac{2x+1}{x+3}=2.
\]

这里的 \(\infty\) 不是某个叫``无穷''的神秘实数点。它是一种简写，表达``当 \(x\) 任意大时，函数值越来越接近某个有限数''。几何上对应水平渐近线 \(y=2\)。

类似地，\(\lim_{x\to -\infty}f(x)\) 研究输入走向负无穷远时的趋势。对于偶函数和奇函数，左右无穷远处的行为可能通过对称性关联起来。

\subsection{图表能提示趋势，不能最终裁决}

你画一个函数值的表格，或者看一眼图像，感觉趋势``很像是收敛到 2''——然后呢？写在卷子上？``看起来像''不是数学论证。

数值采样的危险性在于有限分辨率。函数
\[
f(x)=\sin\frac{1}{x}
\]
在 \(x\to 0\) 时的行为，如果你只算 \(x=0.1, 0.05, 0.02, 0.01\)：
\[
\sin(10)\approx -0.544,\quad \sin(20)\approx 0.913,\quad \sin(50)\approx -0.262,\quad \sin(100)\approx -0.506.
\]
这几个数字跳来跳去，没有规律。实际上，当 \(x\to 0\) 时，\(\sin(1/x)\) 在 \(-1\) 和 \(1\) 之间无限次振荡——无论你靠得多近，振荡永不停止。任何有限的数值表格都不可能揭示这个事实，因为你总是只采样到有限个点。

图像同样不可靠。在屏幕上画 \(\sin(1/x)\)，当 \(x\) 很小时，振荡频率超过像素分辨率，图形被涂成一片蓝色区域。你看到的``填充区域''是人类视觉的错觉，不是函数在收敛。

数值表和图像是探路的罗盘，但极限的存在性需要比``看起来像''更强的保障。你需要的是一种逻辑，能够说：\textbf{从某个足够小的尺度开始，函数值的波动不会超过任意指定的误差界限。}

这就是下一节要展开的严格定义。

\subsection{序列视角：用离散点逼近连续极限}

另一种思考极限的方式是通过序列。如果 \(\lim_{x\to a}f(x)=L\)，那么随便挑一列数 \(x_1, x_2, x_3,\ldots\)，只要每一项都不等于 \(a\) 且整列数收敛到 \(a\)，对应的函数值序列 \(f(x_1), f(x_2), f(x_3),\ldots\) 就必须收敛到 \(L\)。

反过来，如果这个性质对所有这样的序列都成立（在实数范围里），那么函数的极限就是 \(L\)。序列刻画和 \(\varepsilon\)-\(\delta\) 定义是等价的——从不同方向描述同一件事。

序列视角最实用的用途不是证明极限存在，而是证明极限\textbf{不存在}。你只需找到两条通向 \(a\) 的路径，让函数值走向不同的目的地。

对 \(\sin(1/x)\) 在 \(x\to 0\)：
\begin{itemize}
  \item 取 \(x_n = \frac{1}{2\pi n + \pi/2}\)：所有 \(n\) 都满足 \(\sin(1/x_n)=\sin(2\pi n+\pi/2)=1\)。序列收敛到 1。
  \item 取 \(x_n = \frac{1}{2\pi n - \pi/2}\)：所有 \(n\) 都满足 \(\sin(1/x_n)=\sin(2\pi n-\pi/2)=-1\)。序列收敛到 \(-1\)。
\end{itemize}

两条路径都逼近 0，但函数值分别恒等于 1 和 \(-1\)。趋势不一致，极限不存在。这个论证不依赖图像分辨率，不依赖数值表格的采样密度——它依靠的是逻辑上不可避免地暴露了矛盾。

\subsection{这张地图到此为止覆盖了什么}

极限之所以被发明出来，是因为微积分遭遇了一类反复出现的问题：需要知道函数在某个点上的行为，但那个点上函数要么没定义，要么公式失效。核心策略是\emph{不直接看那个点，而是看它附近的所有点，并追问这些点上的函数值是否趋于一致}。

你手中的工具：
\begin{itemize}
  \item 区分函数值与极限值——两者独立，互不绑定。
  \item 单侧极限判断趋势的方向一致性；双侧极限要求两边同意。
  \item 无穷极限描述无界增长（竖直渐近线），极限在无穷远处描述长期行为（水平渐近线）。
  \item 数值表和图像提供线索，不能提供证明。
  \item 序列方法可以严格否定极限的存在。
\end{itemize}

极限概念目前还有一个重大缺口：以上所有论证依赖的都是``越来越接近''、``趋于''、``稳定下来''这些直觉用语。但``接近''到底是什么标准？多近才算``够近''？如果有人说 \(\lim_{x\to 1}x^2=1.999\)，你如何用逻辑而非画图证明他是错的？

直觉阶段到此为止。你需要一套能把``趋近''翻译成精确数学条件的语言。
